%% Appendix sections
%% Note: Sections previously in the appendix (Extended Related Works, Full Experimental Setups,
%% Association Score Computation, Evaluation Metrics, Algorithms, Qualitative Graph Visualization,
%% Structure Learning Examples, and Token Usage Analysis) have been migrated to the main body
%% for the TKDD journal version.


\section{Detailed Prompts for StructSynth}
\label{app:prompts}
This section provides a guide to the prompt templates used in the StructSynth framework. A key step in our methodology is the textualization of diverse data structures---including statistical scores, dependency graphs, and tabular data---to serve as inputs for the Large Language Model (LLM). \hyperref[box:textualization_examples]{Box 1} shows concrete examples of this process. 

The StructSynth pipeline employs a sequence of prompts to complete its tasks. First, the \hyperref[prompt:source]{\texttt{$\pi_\mathtt{source}$}} prompt initiates graph discovery by identifying initial root nodes. During the iterative graph construction, the \hyperref[prompt:generation]{\texttt{$\pi_\mathtt{generation}$}} prompt proposes new dependency edges with corresponding rationales. If a cycle is detected, the \hyperref[prompt:resolve]{\texttt{$\pi_\mathtt{resolve}$}} prompt is used to analyze the conflict and remove the weakest link. Finally, for data synthesis, the \hyperref[prompt:data_gen]{\texttt{$\pi_\mathtt{data\_gen}$}} prompt generates values for dependent features, after which the \hyperref[prompt:data_gen_iso]{\texttt{$\pi_\mathtt{data\_gen\_iso}$}} prompt generates values for the remaining isolated features.


\begin{tcolorbox}[
    float*=t, width=\textwidth, floatplacement=htbp,
    colback=teal!10!white,  % Light teal background
    colframe=teal!60!black, % Dark teal frame
    title={Examples of Data Textualization for LLM Prompts}, % Box title
    fonttitle=\bfseries\color{white}, % Title font bold, color white
    arc=2mm, % rounded corners
    boxrule=1pt % thickness of the frame
] 

\small % Use a smaller font size for the content within the box

\textbf{1. Textual Representation of Statistical Evidence (Association Scores):}
\begin{verbatim}
Association scores for 'Education', scaled 0 to 1.
(Levels: >0.8 Very High, >0.6 High, >0.4 Moderate, >0.2 Low, <=0.2 Very Low)

- Occupation: 0.85 (Very High)
- Salary: 0.72 (High)
- Marital-Status: 0.55 (Moderate)
- Hours-per-week: 0.31 (Low)
- Sex: 0.12 (Very Low)
\end{verbatim}

\vspace{\medskipamount} % Add space

\textbf{2. Textual Representation of a Dependency Graph:}
\begin{verbatim}
Current dependency graph structure with rationales:
- Age -> Marital-Status
Rational: A person's age is a primary factor influencing their marital status.

- Sex -> Relationship: 
Rational: Gender often determines the relationship role within a household.

- Country -> Race: 
Rational: Native country has a strong statistical correlation with a person's race.

- Race -> Education: 
Rational: Statistical data often shows a correlation between a person's race 
and their level of educational attainment.

\end{verbatim}

\vspace{\medskipamount} % Add space

\textbf{3. Textual Representation of Tabular Data (Markdown):}
\begin{verbatim}
| age | education   | occupation      | sex    | hours-per-week | salary |
|-----|-------------|-----------------|--------|----------------|--------|
| 39  | Bachelors   | Adm-clerical    | Male   | 40             | <=50K  |
| 50  | Bachelors   | Exec-managerial | Male   | 13             | <=50K  |
| 28  | Assoc-acdm  | Prof-specialty  | Female | 45             | <=50K  |
\end{verbatim}
\label{box:textualization_examples}
\end{tcolorbox}




\begin{tcolorbox}[
    float*=t, width=\textwidth, floatplacement=htbp,
    colback=yellow!10!white,  % Light yellow background
    colframe=yellow!75!black, % Dark gold/ochre frame
    title={Prompt for Source Nodes Initialization of Tabular Data ($\pi_\mathtt{source}$)}, % Box title
    fonttitle=\bfseries\color{white}, % Title font bold, color black 
    arc=2mm, % rounded corners
    boxrule=1pt % thickness of the frame
]
\small % Use a smaller font size for the content within the box

You are an expert data analyst. Based on the following information:

\vspace{\medskipamount}

\textbf{Task Description:}

Identify the initial "source nodes" for a dependency graph. A source node is a foundational variable that is not caused by any other variable in this dataset.


\vspace{\medskipamount}

\textbf{All Feature Descriptions:}
\begin{verbatim}
{all_features}
\end{verbatim}

\vspace{\medskipamount}

\textbf{Example Data:}
\begin{verbatim}
{few_shot_data}
\end{verbatim}

\vspace{\medskipamount}

\textbf{Your Task is to Identify the Source Nodes:}
\begin{itemize}[leftmargin=*, itemsep=2pt]
    \item Your only task is to analyze the provided column descriptions and data.
    \item Identify the features that are most likely to be \textbf{source nodes} (i.e., fundamental attributes that are not effects of other features).
    \item Provide your answer as a simple list of the identified source node names.
\end{itemize}
\label{prompt:source}
\end{tcolorbox}



\begin{tcolorbox}[
    float*=t, width=\textwidth, floatplacement=htbp,
    colback=pink!10!white,  % Light pink background
    colframe=pink!75!black, % Dark pink/magenta frame
    title={Prompt for Link Generation in Dependency Discovery ($\pi_\mathtt{generation}$)}, % Box title
    fonttitle=\bfseries\color{white}, 
    arc=2mm, % rounded corners
    boxrule=1pt % thickness of the frame
]
\small % Use a smaller font size for the content within the box

You are an expert data analyst building a dependency graph. Based on the following information:

\vspace{\medskipamount} % Add space

\textbf{Task Description:}

For the given $\texttt{\{current\_feature}\}$, your task is to propose its direct successors (effects) 
from the list of available features. For each proposed dependency, you must provide a 
brief, evidence-based rationale.

\vspace{\medskipamount} % Add space

\textbf{Current Graph State:}
\begin{verbatim}
{structure_graph}
\end{verbatim}
\vspace{\medskipamount} % Add space

\textbf{All Candidate Features:}
\begin{verbatim}
{all_features}
\end{verbatim}
\vspace{\medskipamount} % Add space

\textbf{Statistical Evidence (Association Scores):}
\begin{verbatim}
{association_scores}
\end{verbatim}
\vspace{\medskipamount} % Add space

\textbf{Example Data:}
\begin{verbatim}
{few_shot_data}
\end{verbatim}
\vspace{\medskipamount} % Add space

\textbf{Your Task is to Propose and Justify New Links:}
\begin{itemize}[leftmargin=*, itemsep=2pt]
    \item Your focus is on the feature: $\texttt{\{current\_feature}\}$.
    \item Review the \textbf{Statistical Evidence}, which shows the strength of association between $\texttt{\{current\_feature\}}$ and other candidate features $\texttt{\{all\_features\}}$.
    \item Based on the statistical scores, the example data, and the existing graph, identify which other features are most likely to be \textbf{direct effects} of $\texttt{\{current\_feature\}}$.
    \item For each dependency you propose, you \textbf{must} provide a concise $\texttt{rationale}$. This rationale should explain why the dependency makes sense (e.g., "Higher Education is strongly correlated with and typically precedes higher Income").
    \item Do not propose links that would create an obvious logical contradiction with the existing graph structure.
    \item Format your final output as a list of dependency edges and their corresponding rationales.
\end{itemize}

\label{prompt:generation}
\end{tcolorbox}

\begin{tcolorbox}[
    float*=t, width=\textwidth, floatplacement=htbp,
    colback=green!10!white,  % Light green background
    colframe=green!60!black, % Dark green frame
    title={Prompt for Cycle Resolution in Dependency Discovery ($\pi_\mathtt{resolve}$)}, % Box title
    fonttitle=\bfseries\color{white}, % Title font bold, color white
    arc=2mm, % rounded corners
    boxrule=1pt % thickness of the frame
]
\small % Use a smaller font size for the content within the box

You are a logical reasoning expert tasked with resolving a contradiction in a dependency graph. Based on the following information:

\vspace{\medskipamount} % Add space

\textbf{Task Description:}

Analyze the provided dependencies and their rationales within the cycle. Identify the 
single dependency with the weakest or least plausible rationale and recommend it for 
removal to break the cycle.

\vspace{\medskipamount} % Add space

\textbf{Conflicting Cycle with Rationales:}
\begin{verbatim}
{cycle_with_rationales}
\end{verbatim}
\vspace{\medskipamount} 

\textbf{Your Task is to Identify the Weakest Link:}
\begin{itemize}[leftmargin=*, itemsep=2pt]
    \item A logical contradiction (a cycle) has been detected in the graph, detailed in \textbf{Conflicting Cycle with Rationales}.
    \item Your task is to carefully evaluate the \texttt{rationale} provided for each edge in this cycle.
    \item Based on your analysis of the rationales and any relevant context from the column descriptions or data, identify the \textbf{single weakest link}.
    \item The weakest link is the dependency whose rationale is the least plausible, least supported, or most likely to be spurious.
    \item Provide your final answer as the single dependency edge that should be removed to resolve the contradiction.
\end{itemize}

\label{prompt:resolve}
\end{tcolorbox}

\begin{tcolorbox}[
    float*=t, width=\textwidth, floatplacement=htbp,
    colback=blue!5!white,  % Light blue background
    colframe=blue!75!black, % Dark blue frame
    title={Prompt for Structure-Guided Data Synthesis ($\pi_\mathtt{data\_gen}$)}, % Box title
    fonttitle=\bfseries\color{white}, % Title font bold, color white
    arc=2mm, % rounded corners
    boxrule=1pt % thickness of the frame
]
\small 

You are a helpful AI assistant that generates realistic tabular data based on structural dependencies.

\vspace{\medskipamount} % Add space

\textbf{Task Description:}

Generate a realistic synthetic table containing exactly the columns listed in 
$\texttt{\{features\_to\_generate\}}$. This generation must be conditioned on the provided values 
of their parent features.

\vspace{\medskipamount} % Add space

\textbf{Conditioning Data (Current Results):}
\begin{verbatim}
{current_result}
\end{verbatim}
\vspace{\medskipamount} % Add space

\textbf{Relevant Dependency Structure:}
\begin{verbatim}
{related_structure_graph}
\end{verbatim}
\vspace{\medskipamount} % Add space

\textbf{Example Data:}
\begin{verbatim}
{related_few_shot_data}
\end{verbatim}
\vspace{\medskipamount} % Add space

\textbf{Your Task is to Generate the Synthetic Data:}
\begin{itemize}[leftmargin=*, itemsep=2pt]
    \item Your primary task is to generate realistic data for the features listed in $\texttt{\{features\_to\_generate\}}$.
    \item This generation must be \textbf{conditioned} on the data provided in \textbf{Conditioning Data (Current Results)}. These are the values of the parent nodes that influence the features you are generating.
    \item The \textbf{Relevant Dependency Structure} shows you exactly how the conditioning features relate to the features you need to generate.
    \item Use the \textbf{Example Data} to understand the format, range, and statistical properties of the real data.
    \item The generated table should contain columns for exactly the features listed in $\texttt{\{features\_to\_generate\}}$.
    \item Ensure the generated values are plausible and respect the learned dependencies.
\end{itemize}

\label{prompt:data_gen}
\end{tcolorbox}

\begin{tcolorbox}[
    float*=t, width=\textwidth, floatplacement=htbp,
    colback=purple!10!white,  % Light purple background
    colframe=purple!75!black, % Dark purple frame
    title={Prompt for Independent Feature Synthesis ($\pi_\mathtt{data\_gen\_iso}$)}, % Box title
    fonttitle=\bfseries\color{white}, % Title font bold, color white
    arc=2mm, % rounded corners
    boxrule=1pt % thickness of the frame
]
\small % Use a smaller font size for the content within the box

You are a helpful AI assistant that completes tabular data records by generating values for remaining features.

\vspace{\medskipamount} % Add space

\textbf{Task Description:}

The core, structurally-dependent features of a dataset have already been generated. 
Your task is to generate plausible values for the remaining isolated features, ensuring they are statistically consistent with the core features.

\vspace{\medskipamount} % Add space

\textbf{Conditioning Data (Generated Graph-Based Features):}
\begin{verbatim}
{graph_based_values}
\end{verbatim}
\vspace{\medskipamount} % Add space

\textbf{Features to Generate (Isolated Features):}
\begin{verbatim}
{isolated_features_to_generate}
\end{verbatim}
\vspace{\medskipamount} % Add space

\textbf{Example Data:}
\begin{verbatim}
{related_few_shot_data}
\end{verbatim}
\vspace{\medskipamount} % Add space

\textbf{Your Task is to Generate the Independent Values:}
\begin{itemize}[leftmargin=*, itemsep=2pt]
    \item Your primary task is to generate realistic data for the features listed in $\texttt{\{isolated\_features\_to\_generate}\}$.
    \item This generation must be \textbf{conditioned} on the complete set of core features provided in \textbf{Conditioning Data}.
    \item The features you are generating do not have direct parent-child dependencies in the learned graph, but their values should still be plausible and consistent in the context of the entire data record.
    \item Use the \textbf{Example Data} to understand the typical format and statistical properties of these isolated features.
    \item The generated table should contain columns for exactly the features listed in $\texttt{\{isolated\_features\_to\_generate}\}$.
\end{itemize}
\label{prompt:data_gen_iso}
\end{tcolorbox}

\clearpage


